\section{Positive spectral estimates for the exterior packets}
\label{sec:physical-ward-packet-prices}

The exterior Ward argument uses spectral estimates before full rotation
covariance has been proved. The spectral measure for one coordinate axis
therefore cannot be identified with the spectral measure for another axis.
The link between them will instead be the same tested correlation, with its
reflection sign. We develop that link and the estimates it gives.

Write $\epsilon_n=e^{-n}$. A sequence has exponent $\alpha$ if its norm
is at most $C e^{\alpha n}$ for all sufficiently large $n$, for some fixed
$C$. Constants and the starting index may depend on the fixed fields and
test profiles. Finite sums preserve an exponent; products add exponents;
and an estimate with exponent $\alpha$ also holds with any larger exponent.
These rules follow directly from the triangle and product inequalities.

A cluster below means a finite complex linear combination of field words
at fixed, internally distinct positions. A cluster is \emph{on the
$i$-hyperplane} if every point in every word has $i$th coordinate zero.
It has reflection sign $\eta_i\in\{1,-1\}$ if reflection of its points
and labels in that coordinate gives $\eta_i$ times the cluster.
For each axis $i$, use its reflection-positive chamber and its joint
translation spectral measure, as in the heat-vector construction. Denote
the positive measure belonging to the cluster $A$ by $\mu_i^A$ and write
$E=p_i$. It is supported on $E\geq0$. The positive-time weight makes it
locally finite in energy, including on $0\leq E\leq1$; the measure itself
need not have finite total mass.

For completeness, this measure is obtained without an arbitrary boundary
extension. Let $v_u$ be the chamber vector of $A$ translated by $u$ into
the positive half-space, and let $\rho_u$ be its finite joint spectral
measure. The commuting real-time group supplied by
Section~\ref{sec:complex-heat}, together with the three spatial translation
groups, has the joint measure of Section~\ref{sec:spectral-measure}; only
its Euclidean application is used here. Set
$d\mu_i^A=e^{2uE}d\rho_u$. For $v>u$, the semigroup relation
$v_v=e^{-(v-u)H_i}v_u$ gives
$d\rho_v=e^{-2(v-u)E}d\rho_u$, so the definition is independent of $u$.
Two admissible placements are compared through their maximum. On $E\leq1$
the multiplier is at most $e^{2u}$, proving the asserted finite mass there.

The spectral input used here has the following precise form. If all the
points of $A$ lie above $-u$ in coordinate $i$, then, for each $j\ne i$
and $k\geq0$,
\[
 \int e^{-2uE}|p_j|^{2k}\,d\mu_i^A(p)<\infty
\]
when the inequalities at $-u$ are strict. If $A$ is on the $i$-hyperplane,
there is, for each $j,k$, an integer $D$ such that
\begin{equation}\label{eq:packet-polynomial-moment}
 \int(1+E)^{-D}|p_j|^{2k}\,d\mu_i^A(p)<\infty.
\end{equation}
These assertions use only positive-time evolution, spatial translations,
and the polynomial separated-point derivative bounds of the hierarchy.
Here is the passage from those bounds to the second assertion. The Gram
function of the translated positive-time cluster is a finite sum of
separated correlations. Its $2k$th transverse derivative is bounded by
$C t^{-D_0}$ when the reflected halves have distance $t\in(0,1]$.
The positive spectral representation gives
\[
 \int e^{-tE}|p_j|^{2k}\,d\mu_i^A\leq C t^{-D_0}.
\]
To justify the moment before differentiating under the integral, take the
$k$th forward difference of the translated vector. Its squared norm is
the integral of $|e^{ihp_j}-1|^{2k}/|h|^{2k}$ against the damped measure.
The iterated fundamental theorem of calculus bounds this squared norm by
the corresponding $2k$th Gram derivative. Fatou gives the asserted moment
bound. Dominated convergence, using $|e^{ix}-1|\leq |x|$, then gives the
moment identity. Finally integrate the preceding bound against
$t^{D_0}\,dt$ on $(0,1)$. For $E\geq0$,
\[
 \int_0^1t^{D_0}e^{-tE}\,dt
 \geq \frac{e^{-1}}{D_0+1}(1+E)^{-D_0-1}.
\]
Tonelli proves \eqref{eq:packet-polynomial-moment}. For a cluster with a
strict positive-time margin, the same finite-difference argument at that
fixed margin gives the exponentially weighted statement.

\subsection{Moment subtraction with a positive response}

For a real smooth compactly supported profile $f$, put
$f_s(x)=s^{-1}f(x/s)$, $s>0$. Choose once an even nonnegative smooth
function $\varphi$, supported where $1\leq|x|\leq2$, with integral one.
Let $L_j$, $0\leq j<m$, be the Lagrange polynomials for the nodes
$(j+2)^2$. Define a correction $B_{m,s}f$ by
\begin{equation}\label{eq:packet-correction}
 B_{m,s}f=
 \sum_{j<m}\sum_{k\geq0}[X^k]L_j\,
       s^{-2k}\!\left(\int x^{2k}f(x)\,dx\right)
       \bigl(x\mapsto x^{-2k}\varphi(x)\bigr)_{s(j+2)}.
\end{equation}
The second sum is finite. All correction profiles vanish near zero.
Interpolation of the monomials $X^r$, $r<m$, shows that
$R_{m,s}f=f-B_{m,s}f$ has zero moments of orders below $2m$ when $f$ is
even. In fact the even moment of order $2r$ of a term with index $j,k$
is $(s(j+2))^{2r}$ times the corresponding moment of
$x^{-2k}\varphi$, multiplied by its coefficient. Summing the coefficients
is exactly $\sum_j L_j(X)(j+2)^{2r}=X^r$; the odd moments vanish by parity.

For an even real profile let $\mathcal Lf(E)$ be its temporal Laplace
response in the spectral formula. With the Fourier normalization used
there, for $E>0$ it is
\[
 \mathcal L f(E)=2E\int_{\mathbb R}\frac{f(x)}{x^2+E^2}\,dx.
\]
Set $M_{m,s}f=\mathcal L(R_{m,s}f)$. The rational interpolation identity
\[
 \frac1{x+y}-\sum_{j<m}\frac{L_j(x)}{(j+2)^2+y}
 =\frac1{x+y}\prod_{j<m}\frac{(j+2)^2-x}{(j+2)^2+y}
\]
follows by multiplying by $x+y$: the two polynomials agree at the $m$
nodes and at $x=-y$. Scaling and integrating it gives
\begin{equation}\label{eq:packet-positive-kernel}
 M_{m,s}f(E)=2E\int\!\!\int
 \frac{\varphi(v)f(x)}{x^2+E^2}
 \prod_{j<m}\frac{(sv(j+2))^2-x^2}{(sv(j+2))^2+E^2},dx\,dv.
\end{equation}
If $f\geq0$ and its support is in $[-s,s]$, every factor in this product
lies between zero and one. The response is nonnegative. The identity and
its bounds extend to $E=0$ by continuity of the Laplace response; this
retains any spectral mass at zero.

For fixed $m$ and $D\leq2m$, the same product gives, whenever $s\geq1$
and the support of $f_w$ lies in $[-s,s]$,
\begin{equation}\label{eq:packet-response-ceiling}
 0\leq M_{m,s}f_w(E)
 \leq C_{m,f,D}(w^{-1}+s^D)(1+E)^{-D}.
\end{equation}
For $E\leq1$, bound $f_w$ by $w^{-1}\|f\|_\infty$ and use
$\int(x^2+E^2)^{-1}dx=\pi/E$. For $E\geq1$, the product is at most
$C_m\min(1,(s/E)^{2m})$. The remaining integral, after multiplication
by $2E$, is bounded by $2\int f$ because $E\geq1$. These are the two
terms in \eqref{eq:packet-response-ceiling}, with a larger fixed constant.
Combining it with \eqref{eq:packet-polynomial-moment} gives
\begin{equation}\label{eq:packet-response-moment}
 \int M_{m,s}f_w(E)|p_j|^{2k}\,d\mu_i^A
 \leq C(w^{-1}+s^D),
\end{equation}
where $D$ is chosen first and $m$ is then chosen sufficiently large.
In particular $D$ does not grow with the subsequently chosen subtraction
order $m$.
The ordinary double integrals needed for the spectral representation are
also absolutely integrable. Zero moments through order $2m-1$ imply
$|\widehat{R_{m,s}f}(t)|\leq C|t|^{2m}$ near zero by Taylor's formula;
away from zero it is Schwartz. Thus its absolute Laplace integral is
at most $C'(1+E)^{-2m-1}$. Choose $2m+1\geq D$ and apply
\eqref{eq:packet-polynomial-moment}. This justifies Fubini before using
reflection-sign cancellation, rather than inferring it from a signed integral.


There is also a lower bound. Take an even nonnegative profile $\chi$
equal to one on $[-1,1]$ and supported in $[-2,2]$. For $s\geq1$
containing the support of $\chi_w$,
\begin{equation}\label{eq:packet-response-floor}
 M_{m,s}\chi_w(E)\geq
 2^{-m}(E+w)^{-1}(1+E)^{-2m}\quad(E\geq0).
\end{equation}
Indeed each kernel factor is at least
$\tfrac12(1+E)^{-2}$ on the indicated supports. In the remaining integral
restrict $x$ to $[0,wE/(E+w)]$. There $\chi_w=w^{-1}$ and
$x^2+E^2\leq2E^2$, which gives
$2E\int\chi_w(x)/(x^2+E^2)\,dx\geq(E+w)^{-1}$.
Continuity gives the endpoint $E=0$.

\subsection{Changing the spectral axis by its reflection sign}

For $Q\subset\{1,2,3,4\}$, $i\in Q$, and profiles $f_j$, write
\[
 P_i(h;f,Q)=\int\mathcal Lh(p_i)
                   \prod_{j\in Q\setminus\{i\}}f_j(p_j)\,d\mu_i^A(p).
\]
We use this expression only when its absolute integrability has been
established by \eqref{eq:packet-response-moment} or by a positive-time gap.
If $A$ has opposite reflection signs at $i$ and $j$, testing the same
cluster autocorrelation in those two axes gives
\[
 P_i(f_i;f,Q)=-P_j(f_j;f,Q).
\]
For the full set of axes this is the spectral representation of the same
ordinary integral on the two sides; their signs differ by reflection.
For a smaller $Q$, insert in each omitted coordinate the bounded profile
$\chi(p/(N+1))$. It tends pointwise to one, and dominated convergence
with the already integrable temporal response proves the stated identity.
No relation between the two spectral measures has been assumed.

Suppose $f_k\geq0$ and the residual responses of $f_i-b_i$ and $f_j-b_j$
are nonnegative. Apply the identity to those two residuals and expand the
two differences. Dropping the nonnegative full $j$-response gives
\begin{equation}\label{eq:packet-axis-alternation}
 P_i(f_i-b_i;f,Q)\leq
 P_i(f_i-b_i;f[j\leftarrow b_j],Q)
 +P_j(f_j-b_j;f[i\leftarrow b_i],Q).
\end{equation}
A signed correction on the right can be replaced by any pointwise
nonnegative majorant, since all its other factors are nonnegative.

Here is why this inequality buys arbitrarily small losses even if the
original profile widths shrink exponentially. Suppose both reflection
signs occur in $Q$, $A$ is on every hyperplane indexed by $Q$, and
$w_j\leq0$. For every $\delta>0$, sufficiently high subtraction gives
exponent $\delta$ for the residual response in axis $i$, with transverse
profiles $(f_j)_{e^{w_jn}}$.

We spell out the scale order. More generally allow $-K\leq w_j\leq a$,
and designate a set $U\subset Q$ containing every axis whose width may
shrink. Outside $U$ require $w_j\geq0$. There is a positive admissible
correction exponent $a$. For $0\leq d\leq a$, $s\geq1$, require also
$w_i\leq d$ and eventual containment of the scaled $i$-profile in
$[-s e^{dn},s e^{dn}]$. Then this correction scale works after sufficiently
many moments have been subtracted. If $U$ is empty, use \eqref{eq:packet-response-moment}
with $k=0$. The transverse factors are uniformly bounded and the response
has exponent $\max(-w_i,dD)$. Choose $aD<\delta$; the claim follows.

For the induction, change from $i$ to an axis $j$ with the opposite sign
and apply the already established estimate with shrinking set
$U\setminus\{i\}$. Let $a'>0$ be its admissible correction exponent.
Choose first a transverse moment $2k$ so large that
\[
 a'k>K\,|(Q\setminus\{i\})\setminus\{j\}|+\max(K,\delta).
\]
Choose the corresponding moment exponent $D$ in
\eqref{eq:packet-response-moment}, and only then choose
$0<a<a'$ with $aD<\delta$.
The finite formula \eqref{eq:packet-correction} shows that the correction
of $(f_i)_{e^{w_in}}$ at $s e^{dn}$ is dominated in absolute value by
$g_{e^{dn}}$, for one fixed even nonnegative $g$ vanishing near zero:
its coefficients are bounded since $w_i\leq d$, and their supports form
a finite union of fixed annuli. Apply the induction hypothesis at axis
$j$ to this corrected $i$-profile. This bounds the second term of
\eqref{eq:packet-axis-alternation} with exponent $\delta$.

The correction at axis $j$ is bounded by $h_{e^{a'n}}$ with fixed $h$
vanishing near zero. Consequently
\[
 |h_{e^{a'n}}(p_j)|\leq C e^{-2ka'n}|p_j|^{2k}.
\]
The omitted normalization $e^{-a'n}$ only improves this bound because
$a'>0$. Each other transverse factor costs at most $Ce^{Kn}$.
Equation~\eqref{eq:packet-response-moment} now bounds the first term of
\eqref{eq:packet-axis-alternation} with exponent at most
\[
 -2ka'+K\,|(Q\setminus\{i\})\setminus\{j\}|+
       \max(-w_i,dD)<\delta.
\]
This proves the induction step. To arrange the finite induction, visit
an axis of $U$ and then erase it; if its sign is the same as the present
axis, first pass through an opposite-sign axis. Every such intermediate
pass is allowed because both signs occur in $Q$. Each erasure reduces
$|U|$, so the process ends. Moment orders and growing widths have been
chosen in reverse order along this finite route, before $n$ tends to infinity.

Apply this estimate after replacing $f_i$ by $\chi$ and its width by
$e^{qn}$, $q<0$, and take a fixed correction scale ($d=0$).
Equation~\eqref{eq:packet-response-floor}, restricted to $0\leq E\leq1$,
proves the useful point estimate
\begin{equation}\label{eq:packet-point-price}
 \int_{E\leq1}\frac{\prod_{j\in Q\setminus\{i\}}
       (f_j)_{e^{w_jn}}(p_j)}{E+e^{qn}}\,d\mu_i^A(p)
       =O(e^{\delta n}).
\end{equation}
This holds for every $\delta>0$ and all nonnegative even compactly
supported transverse profiles. The constants may depend on their shapes.

\subsection{Keeping a genuine temporal gap}

A collinear cluster is on each hyperplane perpendicular to its line,
but need not be on the hyperplane perpendicular to the line itself.
We cannot use arbitrarily small temporal displacements on that axis.
Fix $u\geq0$ so that translating the cluster by $u$ places all its points
strictly in the positive half-space. Set $c=\varphi*\varphi$ and
$\Phi=\widehat c=(\widehat\varphi)^2$. Since $\varphi$ is real and
even, $\Phi$ is nonnegative, even, rapidly decreasing, and $\Phi(0)=1$.
Choose a smooth even cutoff $\zeta$ equal to one on $[-2u-1,2u+1]$
and zero outside $[-2u-2,2u+2]$. Let $b_w$ be the real profile with
Fourier transform $\zeta(t)\Phi(wt)$, and put $r_w=c_w-b_w$.
Then
\[
 \widehat r_w(t)=(1-\zeta(t))\Phi(wt),\qquad
 \mathcal Lr_w(E)=\int e^{-|t|E}(1-\zeta(t))\Phi(wt)\,dt\geq0.
\]
The transform vanishes on the required gap. Its response is bounded by
$Cw^{-1}e^{-2uE}$. Also $\|b_w\|_\infty\leq C_u$ uniformly in $w$:
Fourier inversion bounds it by a fixed multiple of the $L^1$ norm of
$\zeta(t)\Phi(wt)$, and $\Phi$ is bounded.

For all sufficiently small $w$, uniformly on $0\leq E\leq1$,
\begin{equation}\label{eq:packet-gap-floor}
 \mathcal Lr_w(E)\geq c_u(E+w)^{-1}.
\end{equation}
To see this, put $R=2u+2$ and choose $\gamma>0$ such that
$\Phi(z)\geq1/2$ for $|z|\leq3\gamma$. Take $wR\leq\gamma$ and
integrate over $R<t<R+\gamma/(E+w)$. On this interval $wt\leq2\gamma$,
$Et\leq R+\gamma$, and $\zeta(t)=0$. Thus the integrand is at least
$e^{-(R+\gamma)}/2$, and the interval has length $\gamma/(E+w)$.

Suppose the remaining three axes have both reflection signs and contain
an axis $j$ of sign opposite to the gap axis $i$. The preceding finite
sign-change argument, applied only to those three axes, provides a growing
subtraction scale $e^{a'n}$ at $j$, with any prescribed loss $\delta>0$.
Use \eqref{eq:packet-axis-alternation} with $f_i=c_{e^{qn}}$,
$b_i=b_{e^{qn}}$, and $q<0$. Its second term is bounded by the uniform
constant $\|b_{e^{qn}}\|_\infty$ times that three-axis response; hence it
has exponent $\delta$. Its first term contains the annular correction
$h_{e^{a'n}}(p_j)$. For every fixed $k$,
\[
 \int\mathcal Lr_{e^{qn}}(E)|p_j|^{2k}\,d\mu_i^A
 \leq C e^{-qn}\int e^{-2uE}|p_j|^{2k}\,d\mu_i^A.
\]
The last integral is finite. If the other transverse widths have
$w_l\geq-K$, the first term therefore has exponent at most
$-2ka'+K\,|(\{1,2,3,4\}\setminus\{i\})\setminus\{j\}|-q$.
Choose $k$ large enough to make this at most $\delta$.
Equation~\eqref{eq:packet-gap-floor} gives the point estimate with a
genuine gap in axis $i$. The measures have stayed fixed throughout.

More explicitly, let the cluster lie on a line in direction $l$, have
both reflection signs among the axes other than $l$, and have an
opposite-sign axis to $l$ among those axes. For arbitrary $w_j\leq0$,
$q<0$ and nonnegative even compactly supported $f_j$, the full low-energy
response in axis $i$ has exponent
\begin{equation}\label{eq:packet-line-price}
 \begin{cases}\delta,&i=l,\\ -w_l+\delta,&i\ne l.\end{cases}
\end{equation}
For $i=l$ this is the gap argument. Otherwise apply
\eqref{eq:packet-point-price} with $Q$ the three axes other than $l$,
then bound the remaining factor by
$\|f_l\|_\infty e^{-w_ln}$. In the exterior application the profiles
have shifted exponents $\nu_j=1+w_j$, where $w_l\geq-150$. Thus with
$\delta=1$ the missing-axis exponent is $-\nu_l+1=-w_l\leq150$;
the line-axis exponent is $1$. These bounds hold even when the resolvent
regularizer is $e^{-3000n}$.

\subsection{The exterior of the low spectral box}

This part needs no subtraction or mixed-sign assumption. Fix complex
Schwartz profiles $f_j$, exponents $w_j\geq-K$, and an axis $d$ with
$w_d<0$. Assume $\widehat f_d$ vanishes on $(-1,1)$. Use the temporal
profile $f_d(\,\cdot\,/e^{w_dn})$ and the transverse profiles
$e^{-w_jn}f_j(\,\cdot\,/e^{w_jn})$. If $|b|=1$, their temporal response
$L_n(E)$ satisfies
\[
 |L_n(E)|\leq \|\widehat f_d\|_1\exp(-E/e^{w_dn}),\qquad E\geq0.
\]
Indeed the Fourier transform of the temporal profile is
$e^{w_dn}\widehat f_d(e^{w_dn}t)$, supported where
$|t|\geq e^{-w_dn}$, with unchanged $L^1$ norm; the sign factor $b$
has modulus one.

Put
\[
 C_n=\{p:0\leq p_d\leq1,\quad
                 |p_j|\leq e^{(1+w_j)n}\ (j\ne d)\}.
\]
The product of the three transverse factors is at most $Ce^{3Kn}$.
Outside the indicated transverse box, at least one factor has normalized
argument of size greater than $e^n$. Its Schwartz seminorm of order $r$
therefore improves this to $C_r e^{(3K-r)n}$. Given any real $y$ and
$R\geq0$, choose $r$ large enough. On $E\leq1$ this proves the bound
$C e^{yn}e^{-RE}$ off $C_n$, with a harmless factor $e^R$.
For $E\geq1$, exponential growth of $e^{-w_dn}$ gives, eventually,
\[
 e^{-E/e^{w_dn}}\leq e^{(y-3K)n}e^{-RE}.
\]
For example choose $n$ so that
$e^{-w_dn}\geq R+|y-3K|n+1$ and use $E\geq1$.
Thus the full packet satisfies
\begin{equation}\label{eq:packet-exterior-price}
 |\mathbf1_{C_n^c}(p)L_n(p_d)
       \prod_{j\ne d}e^{-w_jn}f_j(p_j/e^{w_jn})|
 \leq C e^{yn}e^{-Rp_d}.
\end{equation}
Choose $R$ large enough to place the two fixed clusters strictly inside
the positive-time chamber. The exponentially weighted moment statement
then makes the majorant integrable for each of the four positive measures
in their polarization identity. Integrating and adding that finite identity
proves exponent $y$ for the exterior cross response. In particular $y=-101$
is available, without taking a boundary value at zero energy.

The elementary zeroth response used inside $C_n$ is also explicit. If
$1+w_j\geq-K$, then
\[
 \int_{E\leq1}\frac{\prod_{j\ne d}\chi_{e^{(1+w_j)n}}(p_j)}{E+1}
                  \,d\mu_d^A(p)=O(e^{3Kn}).
\]
Each of three profiles is bounded by $\|\chi\|_\infty e^{Kn}$ and the
remaining integral is finite. With $K=300$ the exponent is $900$.

\subsection{The diagonal-stress estimate and its algebraic reduction}

For every $\delta>0$, even nonnegative smooth compactly supported
transverse profiles $f_j$, and $w_j\leq a\leq0$, the packet application
also needs the improved diagonal-stress bound
\begin{equation}\label{eq:packet-cone-required}
 \int_{E\leq1}\frac{E e^{4an}}{(E+e^{an})^4}
             \prod_{j\ne d}(f_j)_{e^{w_jn}}(p_j)\,d\mu_d^{J_{dd}}(p)
       =O(e^{(\delta+4a)n}),\qquad w_j\leq a\leq0.
\end{equation}
It is a separate consequence of stress conservation, not of
\eqref{eq:packet-point-price}. The following finite algebra explains the
normalization in its proof. Put $w=e^{an}$ and keep one sufficiently large
subtraction order $m$ fixed. The subtracted differences of profiles at
widths $w$ and $2w$ give three real responses $F,G,H$ with
$F\geq0$, $G\leq0$, $H\geq0$. Conservation and the fixed-band estimates
are required to supply real sequences $A,B,V,W,R,T$ such that, with
$x=\delta+2a$, each of
\[
 V-A,\quad W-4B,\quad R-V,\quad T-4W,\quad
 F-(A-B),\quad G-(V-W),\quad H-(R-T)
\]
has exponent $x$. These are estimates on the same subtracted profiles.
They imply $F=O(e^{xn})$: if $E_n$ is the sum of their absolute errors,
then $G\leq0$ gives $A-4B\leq3E_n$ and $H\geq0$ gives
$A-16B\geq-7E_n$. Hence $B\leq(10/12)E_n$ and
$0\leq F\leq A-B+E_n\leq(13/2)E_n$.
The constants can be enlarged freely; the exponent is unchanged.

The positive response kernel of this difference has a lower bound
$c_m E w^2/(E+w)^4$ on $0\leq E\leq1$. This follows directly from
\eqref{eq:packet-positive-kernel}, at the fixed subtraction scale $s=8$.
Write $\varphi_k(x)=x^{-2k}\varphi(x)$ and compare its widths $w$ and
$2w$, with $0<w\leq1$. After scaling back to $x\in\operatorname{supp}
\varphi$, the factor $(x^2w^2+E^2)^{-1}$ and every numerator of the
product decrease when $w$ is replaced by $2w$. The denominators of the
product do not change. Thus every such dyadic response is nonnegative.
For $k=0$ and $E\leq1$, the product at $2wx$ is at least
$2^{-m}2^{-2m}$. Furthermore, since $1\leq|x|\leq2$,
\[
 \frac1{w^2x^2+E^2}-\frac1{4w^2x^2+E^2}
 \geq\frac{w^2}{400(E+w)^4}.
\]
Indeed the numerator of the difference is $3w^2x^2\geq w^2$, and
its positive denominator is at most $400(E+w)^4$.
Integrate against $\varphi(x)\varphi(v)$, both of integral one, and
multiply by $2E$. The bound follows with
$c_m=2\cdot2^{-3m}/400$; $E=0$ follows by continuity.
Multiplying its integral bound
by $c_m^{-1}w^2$ would give exactly
\eqref{eq:packet-cone-required}, with exponent $\delta+4a$.
For the intended choice $a=-99$, $\delta=1$, this is $-395$.
Section~\ref{sec:physical-ward-stress-cone} proves the conservation and
fixed-band estimates that produce these seven errors for the same stress
profiles. Together with the reduction above, it supplies the stated cone
bound used in the exterior packet argument.
